\documentclass{situationist-preprint}

\title{Mathematical Sublimity}
\msid{AES-10-b}
\msfamily{AES-10}
\msdiscipline{Aesthetics}
\mstype{Companion manuscript}
\mskeywords{sublime; mathematics; infinity; large numbers; Kant}

\begin{document}
\maketitle

\begin{abstract}
AES-10-a argues that the sublime remains possible, because its conditions require a magnitude that defeats comprehension while thought grasps it as a determinate whole, and nothing in these conditions requires ignorance of measure. We apply the account to mathematics itself. Very large finite numbers, defined by recursive constructions, and the hierarchy of infinite cardinals present magnitudes whose comprehension fails at once while their structure is grasped exactly. We argue that they are among the clearest cases of the sublime the account allows, that they meet the condition of an attempted comprehension, and that what distinguishes their sublimity from mere bewilderment is the determinacy that proof supplies.
\end{abstract}

\section{The cases}

On Definition~2.2 of AES-10-a, an experience is sublime if it presents an object as having a magnitude, the attempt to comprehend the magnitude fails, the subject grasps it as a determinate whole in thought, and the felt failure is accompanied by pleasure in the grasp. Proposition~3.1 of AES-10-a states that knowledge of measure enables rather than defeats the experience. Kant's mathematical sublime already involves number: the imagination runs through units and fails to hold them together, while reason demands the totality \citep[§§25--27]{Kant1790}.

Consider a number defined by iterated exponentiation, such as a tower of tens of a given height. The definition is short and its value is exact. The attempt to comprehend it begins with apprehension: one writes out the first stage, then the second, and sees that the third cannot be written in any physical medium. Comprehension fails in the precise sense of Definition~2.1 of AES-10-a. Yet the number is grasped as a determinate whole: its definition fixes it, and one can prove facts about it, for instance that it exceeds any number that could be written out digit by digit.

Consider next Cantor's theorem that every set is strictly smaller than the set of its subsets \citep{Cantor1891}. The diagonal argument is short and can be followed step by step. It yields an unending hierarchy of infinite sizes. Here the attempt at comprehension fails twice: the imagination cannot hold even the first infinite totality together, and the theorem shows a larger one beyond it, and another beyond that. Thought nonetheless grasps each as determinate, and grasps the law of the hierarchy.

\begin{proposition}\label{p:math}
Mathematical magnitudes, finite and infinite, can be objects of sublime experience in the sense of Definition~2.2 of AES-10-a.
\end{proposition}

\begin{proof}[Argument]
Conditions (i)--(iii) are met in the cases described, as shown. Condition (iv) is met when the subject takes pleasure in the reach of the definition or proof that grasps what imagination cannot hold. Such pleasure is widely reported by mathematicians and students, and nothing in its object makes it unfitting.
\end{proof}

This supports Theorem~3.3 of AES-10-a, of which the mathematical cases are the limit: informational magnitudes are sizes of spaces of possibilities, and these spaces are described mathematically.

\section{Objections}

\begin{objection}
Mathematical objects are abstract; nothing about them is presented to sense, so condition (ii), which requires an attempt at comprehension in imagination, cannot be met.
\end{objection}

\begin{reply}
The reply of AES-10-a to the objection from presentation applies: the attempt requires only apprehension, which can be successive and symbolic. Writing out the stages of a tower of exponents, or listing the first rows of a diagonal construction, is apprehension through a sensible unit; it is the failure to hold these together that the definition requires.
\end{reply}

\begin{objection}
The experience of proof is intellectual satisfaction, not the sublime.
\end{objection}

\begin{reply}
Satisfaction in a proof alone is not sublime, and Definition~2.2 does not count it so. The sublime requires the felt failure of comprehension as well, and the pleasure in its contrast with the grasp. A proof of a small identity involves no such failure. Cantor's theorem does, and its characteristic response, which \citet{Moore2001} describes in the history of reflection on the infinite, combines vertigo and clarity in the way the definition describes.
\end{reply}

\section{Conclusion}

Mathematics supplies magnitudes that defeat comprehension as decisively as any, together with the most determinate grasp of them available. On the account of AES-10-a, it is therefore a domain in which the sublime is not only possible but exemplary.

\begin{thebibliography}{9}\small
\bibitem[Cantor(1891)]{Cantor1891} Cantor, G. (1891). \"Uber eine elementare Frage der Mannigfaltigkeitslehre. \emph{Jahresbericht der Deutschen Mathematiker-Vereinigung}, 1, 75--78.
\bibitem[Kant(1790)]{Kant1790} Kant, I. (1790). \emph{Critique of the Power of Judgment}. Trans. P. Guyer and E. Matthews. Cambridge University Press, 2000.
\bibitem[Moore(2001)]{Moore2001} Moore, A.~W. (2001). \emph{The Infinite}. 2nd ed. Routledge.
\end{thebibliography}

\end{document}
